% !TEX root = ../algebra_02.tex
\section{Центр конечной $p$-группы. Группы порядка $p^2$}

\subsection*{Центр группы}

\begin{Definition}{Центр группы}{}
	\textbf{Центром} группы $G$ называется множество элементов, коммутирующих со всеми элементами группы:
	\[ Z(G) = \{h \in G \mid \forall g \in G \colon gh = hg\} \] 
	Очевидно, что $Z(G)$ является нормальной подгруппой $G$ ($Z(G) \triangleleft G$).
\end{Definition}

\begin{Theorem}{О центре конечной $p$-группы}{}
	Пусть $G$ — конечная $p$-группа, то есть $|G| = p^n$, где $p$ — простое число, $n \in \mathbb{N}$. Тогда центр группы нетривиален: $Z(G) \neq \{e\}$.
\end{Theorem}

\begin{proof}
	Рассмотрим действие $G$ на себе сопряжениями: $g \cdot m = gmg^{-1}$.
	Орбитами этого действия являются классы сопряженности.
	По следствию из теоремы об орбите и стабилизаторе, размер каждой орбиты $|Gm|$ делит $|G|$, поэтому $|Gm| = p^k$ для некоторого $0 \le k \le n$.
	Разбивая множество $G$ на орбиты, получаем уравнение классов:
	\[ |G| = \sum_{k=0}^{n} m_k p^k = p^n \]
	где $m_k$ — количество орбит размера $p^k$.
	Из этого равенства следует, что $p$ делит $m_0$ (количество орбит длины 1).
	Орбита элемента $m$ имеет длину 1 тогда и только тогда, когда $\forall g \in G \colon gmg^{-1} = m \iff gm = mg$, что означает $m \in Z(G)$.
	Следовательно, $m_0 = |Z(G)|$. Так как нейтральный элемент $e$ всегда лежит в $Z(G)$, имеем $|Z(G)| \ge 1$.
	Но так как $p$ делит $|Z(G)|$, мы получаем, что $|Z(G)| \ge p > 1$. Значит, $Z(G) \neq \{e\}$.
\end{proof}

\begin{Lemma}{Факторгруппа по центру}{}
	Если группа $G$ не является абелевой, то факторгруппа $G/Z(G)$ не может быть циклической.
\end{Lemma}

\begin{proof}
	Докажем от противного.
	Пусть $G/Z(G)$ — циклическая группа, порожденная смежным классом $g_0 Z(G)$:
	\[ G/Z(G) = \langle g_0 Z(G) \rangle \] 
	Тогда любой элемент $g \in G$ представим в виде $g \in g_0^i Z(G)$ для некоторого целого $i$, то есть $g = g_0^i h$, где $h \in Z(G)$.
	Возьмем произвольные два элемента $g, g' \in G$. Они представимы как $g = g_0^i h$ и $g' = g_0^j h'$, где $h, h' \in Z(G)$.
	Рассмотрим их произведение:
	\[ gg' = (g_0^i h)(g_0^j h') = g_0^i g_0^j h h' = g_0^{i+j} h h' \] 
	Мы смогли переставить элементы, так как $h$ и $h'$ лежат в центре и коммутируют со всем.
	Аналогично:
	\[ g'g = (g_0^j h')(g_0^i h) = g_0^j g_0^i h' h = g_0^{i+j} h h' \] 
	Получили $gg' = g'g$ для любых $g, g' \in G$.
	Следовательно, $G$ — абелева группа. Противоречие с условием.
\end{proof}

\subsection*{Группы порядка $p^2$}

\begin{Corollary}{Группы порядка $p^2$}{}
	Если $p$ — простое число, то любая группа $G$ порядка $p^2$ является абелевой.
\end{Corollary}

\begin{proof}
	Так как $G$ — $p$-группа, её центр нетривиален, следовательно, возможные порядки центра: $|Z(G)| = p$ или $|Z(G)| = p^2$.
	Предположим, $|Z(G)| = p$. Тогда индекс подгруппы $(G : Z(G)) = |G/Z(G)| = p$.
	Группа простого порядка всегда циклическая, значит $G/Z(G)$ циклическая. Но по предыдущей лемме, это влечет абелевость $G$.
	Если $G$ абелева, то она совпадает со своим центром, то есть $|Z(G)| = p^2$. Мы получили противоречие с допущением $|Z(G)| = p$.
	Следовательно, единственный возможный случай — это $|Z(G)| = p^2$, то есть $G = Z(G)$, и группа абелева.
\end{proof}
